\documentclass{ximera}
\title{Functions of Two Independent Variables, Introduction}

\newcommand{\pskip}{\vskip 0.1 in}

\begin{document}
\begin{abstract}
Introduction to traces, level curves and surfaces.
\end{abstract}
\maketitle



\begin{example}  \label{Edfhf45}
\begin{enumerate}
\item Move the slider $k$ to see the level curves 
\[
    f(x,y)=k
\]
of the function $z=f(x,y)$. Use the level curves to sketch the surface $z=f(x,y)$. Then check the \emph{Surface} folder in Line 8 to see how you did. Change the setting to \emph{Transparent Surfaces} by clicking the wrench in the upper right corner of the worksheet and selecting \emph{More Options}.


\pdfOnly{
Access Desmos interactives through the online version of this text at
 
\href{https://www.desmos.com/3d/u5syjyxcsw}.
}
 
\begin{onlineOnly}
    \begin{center}
\desmosThreeD{u5syjyxcsw}{900}{600}
\end{center}
\end{onlineOnly}

\href{https://www.desmos.com/3d/u5syjyxcsw}{163: Hemisphere}

\item Deactivate the folder \emph{Surface} in Line 8. Then use the level curves of the function $z = f(x,y)$ above to approximate the partial derivative
\[
    \frac{\partial z}{\partial x} \Big|_{(x,y)= (-2,1)} .
\]
Keep in mind the spacing of the level curves (Line 6) in computing the approximation.

\item Drag the slider $\alpha$ in Line 7 to approximate the partial derivative
\[
    \frac{\partial z}{\partial y} \Big|_{(x,y)= (-2,1)} .
\]

\item Use the expression
\[
  z = f(x,y) = \sqrt{9-x^2}
\]
for the function $f$ to answer the following questions.

\begin{enumerate}
\item Write the domain and range of the function $z=f(x,y)$ in set-builder notation.

\item Evaluate the partial derivatives 
\[
    \frac{\partial z}{\partial x} \Big|_{(x,y)= (-2,1)} 
\]
and
\[
    \frac{\partial z}{\partial y} \Big|_{(x,y)= (-2,1)} 
\]
and compare these with your estimates.
\end{enumerate}

\item Sketch the cross-section of the surface $z=f(x,y)$ cut out by the plane $y=1$. Then use geometry and also the methods of Calculus 1 to find the slope of the tangent to the cross-section at the point $(-2, 1, 2)$. 

Then activate the folders \emph{Surface} (Line 8) and \emph{Trace} (Line 10) to see how you did. But first drag the slider $\alpha$ in Line 7 back to $\alpha = 0$.

\item Sketch the cross-section of the surface $z=f(x,y)$ cut out by the plane $x=-2$. Then use geometry and also the methods of Calculus 1 to find the slope of the tangent to the cross-section at the point $(-2, 1, 2)$. 

Drag the slider $\alpha$ in Line 7 to $\alpha = pi/2$ to see how you did.

\end{enumerate}
\end{example}



\begin{example}  \label{Ed4fgvvvvf45}
Move the slider $k$ to see the level curves 
\[
    f(x,y)=k
\]
of the function $z=f(x,y)$. Use the level curves to imagine or draw the surface $z=f(x,y)$. Then check the \emph{Surface} folder in Line 8 to see how you did. Change the setting to \emph{Transparent Surfaces} by clicking the wrench in the upper right corner of the worksheet and selecting \emph{More Options}.


\pdfOnly{
Access Desmos interactives through the online version of this text at
 
\href{https://www.desmos.com/3d/u5syjyxcsw}.
}
 
\begin{onlineOnly}
    \begin{center}
\desmosThreeD{yj0kpv4zmi}{900}{600}
\end{center}
\end{onlineOnly}


 
\href{https://www.desmos.com/3d/yj0kpv4zmi}{163: Hyperboloid}
\end{example}



\begin{example}  \label{Ed44mlk45}
Move the slider $k$ to see the level curves 
\[
    f(x,y)=k
\]
of the function $z=f(x,y)$. Use the level curves to imagine or draw the surface $z=f(x,y)$. Then check the \emph{Surface} folder in Line 8 to see how you did. Change the setting to \emph{Transparent Surfaces} by clicking the wrench in the upper right corner of the worksheet and selecting \emph{More Options}.


\pdfOnly{
Access Desmos interactives through the online version of this text at
 
\href{https://www.desmos.com/3d/u5syjyxcsw}.
}
 
\begin{onlineOnly}
    \begin{center}
\desmosThreeD{vgyxfvgsyp}{900}{600}
\end{center}
\end{onlineOnly}


 
\href{https://www.desmos.com/3d/vgyxfvgsyp}{163: Cone}
\end{example}




\end{document}